\label{chap:p3-83-26-raiz_menos_uno}

\section*{La racine interne de \texorpdfstring{\(-1\)}{-1} et la famille exponentielle tritique}

La racine de \(-1\) n’est pas introduite comme un symbole formel sans antécédent. Les
six unités modulo neuf sont munies de coordonnées sur
\(\mathbb P^1(\mathbb F_5)\), et le caractère quadratique de
\(\mathbb F_5\) détermine la matrice de conférence de Paley. Sa réduction
ternaire est
\[
 A_W=
 \begin{pmatrix}
 0&1&1&1&1&1\\
 1&0&1&2&2&1\\
 1&1&0&1&2&2\\
 1&2&1&0&1&2\\
 1&2&2&1&0&1\\
 1&1&2&2&1&0
 \end{pmatrix}\in M_6(\mathbb F_3).
\]
L’identité de conférence se réduit à
\begin{equation}\label{p3-83-c26-s1:eq:constantes-raiz-interna-menos-uno}
 \boxed{A_W^2=2I_6=-I_6.}
\end{equation}
Par conséquent, \(A_W\) est un opérateur de quart de tour dans l’espace ternaire.
Après extension des scalaires à \(\mathbb F_9\), il se décompose en les espaces
propres conjugués des racines de \(X^2+1\). La source entière commune
\[
 \mathcal Q_{\mathbb Z}=\mathbb Z[u]/(u^2+1)
\]
possède une réalisation finie \(u\mapsto A_W\) et une réalisation réelle
\(u\mapsto J_{\rm e}\). La source coordonne la relation quadratique ; elle n’identifie pas
des corps de caractéristiques distinctes.

Les trois régimes de TRIT se réalisent par des opérateurs
\[
 J_{+1}^2=-I,\qquad J_0^2=0,\qquad J_{-1}^2=I.
\]
En séparant les puissances paires et impaires de l’exponentielle, on obtient une
seule famille :
\begin{equation}\label{p3-83-c26-s1:eq:constantes-euler-extendido}
 \boxed{
 \exp(tJ_\tau)=C_\tau(t)I+S_\tau(t)J_\tau,
 \qquad \tau\in\{+1,0,-1\},}
\end{equation}
où
\[
 (C_{+1},S_{+1})=(\cos,\sin),\qquad
 (C_0,S_0)=(1,t),\qquad
 (C_{-1},S_{-1})=(\cosh,\sinh).
\]
Ses trois spécialisations directrices sont
\begin{equation}\label{p3-83-c26-s1:eq:constantes-euler-tres-ramas}
 \exp(\pi J_{+1})+I=0,\qquad
 \exp(J_0)=I+J_0,\qquad
 \exp\!\bigl(2(\log\varphi)J_{-1}\bigr)=F_{\rm av}^{\,2}.
\end{equation}
La branche elliptique contient la clôture circulaire complexe ; la parabolique, le
transport de seuil ; l’hyperbolique, l’auto-échelle dorée. L’équation
d’Euler apparaît ainsi comme une spécialisation d’une famille qui coordonne rotation,
frontière et expansion hyperbolique.


\section{Réalisation elliptique de \(J^2=-I\) par la fibration de Hopf et l’holonomie de Möbius}

\subsection{Structure complexe et fibration de Hopf}

Soit \(S^3\subset\mathbb C^2\) et
\[
h(z_0,z_1)=
\bigl(2\Re(z_0\overline z_1),
2\Im(z_0\overline z_1),|z_0|^2-|z_1|^2\bigr)
\in S^2.
\]
Dans le tore de Hopf
\[
T_\eta=\{(z_0,z_1):|z_0|=\cos\eta,
|z_1|=\sin\eta\},
\]
les courbes
\[
\gamma_+(t)=(\cos\eta\,e^{it},\sin\eta\,e^{it}),
\]
\[
\gamma_-(t)=(\cos\eta\,e^{it},\sin\eta\,e^{-it})
\]
présentent des comportements distincts : \(\gamma_+\) est une fibre de Hopf et
\(\gamma_-\) se projette avec un degré deux. Sous projection stéréographique, les
classes homologiques \((1,1)\) et \((1,-1)\) produisent les deux familles diagonales
de Villarceau.

La conjugaison complexe totale
\[
K(z_0,z_1)=(\overline z_0,\overline z_1)
\]
préserve chaque famille non orientée et inverse son orientation. Elle n’échange pas
les deux familles. La conjugaison partielle peut les échanger, mais ce n’est ni une
application complexe linéaire ni l’antiunitaire standard global. Cette correction
est essentielle pour ne pas fabriquer une symétrie particule--antiparticule à partir du
dessin toroïdal.

\subsection{Revêtement double et retour spinoriel}

L’application \(SU(2)\to SO(3)\) est un revêtement double. Un tour de
\(2\pi\) dans \(SO(3)\) peut se relever en \(-I\) dans \(SU(2)\) ; le tour de
\(4\pi\) récupère \(I\). La structure de signe ne transforme pas un cercle en
ruban de Möbius. Pour obtenir un ruban, il faut un fibré réel en droites
dont le caractère d’holonomie soit \(-1\).

Soit \(L_\chi\to S^1\) le fibré associé à
\[
\chi:\pi_1(S^1)\longrightarrow\{\pm1\},
\qquad \chi(1)=-1.
\]
Son espace total est un ruban de Möbius. La route \(C_M=-\operatorname{rev}\)
apporte le candidat de signe ; le cercle de Villarceau apporte le lacet ;
l’entrelaceur physique doit prouver que le premier est l’holonomie du second.

\subsection{Réalisation toroïdale de l’opérateur \texorpdfstring{\(B_{\mathrm{mem}}^{\varepsilon}\)}{Bmem} et détermination de l’angle de Villarceau}

Avec \(B_c=7I+2S\), définissons l’opérateur orienté de mémoire
\[
B_{\rm mem}^{\varepsilon}=5I+2\varepsilon S,
\qquad\varepsilon=\pm1.
\]
Pour \(\varepsilon=+1\),
\[
\Spec(B_{\rm mem}^{+})=\{7,3\},
\qquad\det B_{\rm mem}^{+}=21.
\]
Interpréter les valeurs \((7,3)\) comme les bords équatoriaux extérieur et intérieur
dans une carte positive commune produit
\[
R+r=7\ell,\qquad R-r=3\ell,
\]
et donc
\[
R=5\ell,\qquad r=2\ell,\qquad
\sin\beta=\frac rR=\frac25.
\]
L’angle de Villarceau est
\[
\beta=\arcsin\frac25
=23.5781784782\ldots^\circ.
\]

Le théorème algébrique \(\{7,3\}\leftrightarrow(5,2)\) est exact. La prémisse
géométrique qui interprète le spectre comme les deux bords d’un tore est une
interface déclarée. Elle ne doit pas être dissimulée et l’angle ne doit pas être employé comme cible.

La réalisation du bloc complet \(B_c\) conserve une autre lecture :
\(R:r=7:2\) et \(\beta=\arcsin(2/7)\). Ce ne sont pas des valeurs rivales ; elles relèvent de
deux foncteurs : bloc total et mode de mémoire. La priorité physique du second
doit être décidée par l’entrelaceur avec la dynamique, et non par une proximité décimale avec
une autre constante.

\subsection{Action de l’opérateur \texorpdfstring{\(J^2=-I\)}{J2=-I} sur les coordonnées angulaires conjuguées}

L’angle natif de HMT est une phase \(u\in\mathbb R/\mathbb Z\). Les lecteurs
publient
\[
\theta_{\rm rad}=2\pi u,
\qquad
\theta_{\rm deg}=360u.
\]
Radians et degrés sont des cartes distinctes du même élément circulaire. Le radian
n’est pas plus « ontologique » parce qu’il est adimensionnel ; sa naturalité provient de la
longueur d’arc. Les degrés fournissent une autre coordonnée absolue du
tour. Toute relation avec \(\alpha^{-1}\), \(\varphi\) ou CKM doit
d’abord être formulée comme un quotient ou une holonomie invariante, puis seulement
être évaluée dans une carte angulaire.

L’identité exacte déjà disponible
\[
\tanh^2(2\log\varphi)=\frac59
\]
relie le canal Perron--Fibonacci au quotient spectral de \(B_c\). Son
complément est \(4/9\), le saut normalisé. Cette relation est structurelle.
Elle n’implique pas que tout angle construit avec \(\varphi\) soit un angle physique.

\subsection{Brouwer, Möbius et degré}

Le théorème du point fixe de Brouwer, les transformations fractionnaires de
Möbius et le ruban de Möbius sont trois objets distincts. Ils ne sont reliés
qu’après fixation du domaine, de l’action et du fibré. Les orbites elliptiques d’avance et de
retard peuvent être modélisées comme deux sections d’un fibré à holonomie de
signe ; la transformation de Möbius décrit l’action conforme dans la carte ;
le théorème de Brouwer garantit un point fixe pour les applications continues du disque.
Aucune de ces phrases ne remplace le diagramme de réalisation.

\begin{remark}[Résultat et frontière]
On démontre le degré double de la courbe de Hopf, l’orientation de la
conjugaison totale et la construction algébrique \(7,3\mapsto5,2\). La lecture
toroïdale (5{:}2), l’holonomie de Möbius et l’interprétation
particule--antiparticule forment une chaîne d’interfaces qui doit partager un
unique entrelaceur. Le volume n’impose pas cet entrelaceur au moyen d’une
coïncidence angulaire.
\end{remark}
